\documentclass[12pt]{amsart}

\usepackage[margin=1.1in]{geometry}
\usepackage[T1]{fontenc}
\usepackage{lmodern}
\usepackage{amsmath,amssymb,amsthm,mathtools}
\usepackage{enumitem}
\usepackage[colorlinks=true,linkcolor=blue,citecolor=blue,urlcolor=blue]{hyperref}
\usepackage{microtype}

\newtheorem{theorem}{Theorem}
\newtheorem{conjecture}{Conjecture}
\newtheorem{proposition}{Proposition}
\newtheorem{lemma}{Lemma}
\newtheorem{corollary}{Corollary}
\newtheorem*{remark}{Remark}
\newtheorem*{example}{Example}

\DeclareMathOperator{\spol}{sp}
\DeclareMathOperator{\num}{num}
\DeclareMathOperator{\den}{den}
\DeclareMathOperator{\rad}{rad}
\newcommand{\floor}[1]{\left\lfloor #1\right\rfloor}
\newcommand{\bbfm}{Ballantine--Beck--Feigon--Maurischat}
\newcommand{\ii}{\mathrm{i}}
\newcommand{\OO}{\mathcal O}
\newcommand{\TT}{\mathcal T}
\newcommand{\BB}{\mathcal B}
\newcommand{\ZZ}{\mathbb Z}

\begin{document}

Let \(\lambda=(\lambda_1,\ldots,\lambda_k)\) be a partition of \(n\).  The
\emph{subsum polynomial} of \(\lambda\) is
\begin{equation}
   \spol(\lambda,x):=\prod_{j=1}^k(1+x^{\lambda_j}).
\end{equation}
Equivalently, if \(m_\lambda(i)\) denotes the multiplicity of the part \(i\) in
\(\lambda\), then we have
\[
   \spol(\lambda,x)=\prod_{i\ge1}(1+x^i)^{m_\lambda(i)}.
\]

A common denominator is
$$
   \den^*(n,x):=\prod_{i\ge1}(1+x^i)^{\floor{n/i}},
$$
and the corresponding numerator is
$$
   \num^*(n,x):=\sum_{\lambda\vdash n}h_\lambda^{(n)}(x),
$$
where
$$
   h_\lambda^{(n)}(x):=
   \prod_{i\ge1}(1+x^i)^{\floor{n/i}-m_\lambda(i)}.
$$
One can remove the common divisor of the individual summands 
\[
   G(n,x):=\gcd\{h_\lambda^{(n)}(x):\lambda\vdash n\},
\]
and in turn define
\begin{equation}
   \num(n,x):=\frac{\num^*(n,x)}{G(n,x)},
   \qquad
   \den(n,x):=\frac{\den^*(n,x)}{G(n,x)}.
\end{equation}
We also use the convention \(\num(0,x)=1\) for the empty partition.

We stress that \(G(n,x)\) is not defined to be
\(\gcd(\num^*(n,x),\den^*(n,x))\).  It is the common factor of the summands
\(h_\lambda^{(n)}(x)\).

\begin{conjecture}\label{conj:conj2}  For \(n\ge1\), we have
\[
   \gcd(\num(n,x),\den(n,x))=1.
\]
\end{conjecture}

\end{document}
